\documentclass[12pt,letterpaper]{exam}
\usepackage[lmargin=1in,rmargin=1in,tmargin=1in,bmargin=1in]{geometry}
\usepackage{../style/exams}

% -------------------
% Course & Exam Information
% -------------------
\newcommand{\course}{MATH 142: Exam 3}
\newcommand{\term}{Spring ---\textsubscript{\textsubscript{2}} 2026}
\newcommand{\examdate}{04/16/2026}
\newcommand{\timelimit}{75 Minutes}

\setbool{hideans}{true} % Student: True; Instructor: False

\usepackage{cancel}

\usepackage{cancel}
\usepackage{array}
\newcolumntype{C}[1]{>{\centering\let\newline\\\arraybackslash\hspace{0pt}}m{#1}}

% -------------------
% Content
% -------------------
\begin{document}

\examtitle
\instructions{Write your name on the appropriate line on the exam cover sheet. This exam contains \numpages\ pages (including this cover page) and \numquestions\ questions. Check that you have every page of the exam. Answer the questions in the spaces provided on the question sheets. Be sure to answer every part of each question and show all your work and fully justify all your responses. If you run out of room for an answer, continue on the back of the page --- being sure to indicate the problem number.
} 
\scores
\bottomline
\newpage


% -------------------
% Questions
% -------------------
\begin{questions}

% Question 1
\newpage
\question[20] Fill out the table of Maclaurin series below.
	\begin{table}[!ht]
	\hspace{-0.6cm}
	\begin{tabular}{C{1.71cm}C{0.2cm}C{5cm}C{0.2cm}C{5cm}C{0.2cm}C{2.8cm}}
	{\bfseries Function} && {\bfseries Series} && {\bfseries First Two Nonzero Terms} && {\bfseries Interval of Convergence} \\ \hline
	\\[1.5cm]
	$\dfrac{1}{1 - x}$ && \psol{$\ds\sum_{n=0}^\infty x^n$} && \psol{$1 + x$} && \psol{$(-1, 1)$} \\ \cline{3-3} \cline{5-5} \cline{7-7}
	\\[1.5cm]
	$e^x$ && \psol{$\ds\sum_{n=0}^\infty \dfrac{x^n}{n!}$} && \psol{$1 + x$} && \psol{$(-\infty, \infty)$} \\ \cline{3-3} \cline{5-5} \cline{7-7}
	\\[1.5cm]
	$\sin(x)$ && \psol{$\ds\sum_{n=0}^\infty (-1)^n\, \dfrac{x^{2n + 1}}{(2n + 1)!}$} && \psol{$x - \dfrac{x^3}{3!}$} && \psol{$(-\infty, \infty)$} \\ \cline{3-3} \cline{5-5} \cline{7-7}
	\\[1.5cm]
	$\cos(x)$ && \psol{$\ds\sum_{n=0}^\infty (-1)^n\, \dfrac{x^{2n}}{(2n)!}$} && \psol{$1 - \dfrac{x^2}{2!}$} && \psol{$(-\infty, \infty)$} \\ \cline{3-3} \cline{5-5} \cline{7-7}
	\\[1.5cm]
	$\ln(1 + x)$ && \psol{$\ds\sum_{n=1}^\infty (-1)^{n+1}\, \dfrac{x^n}{n}$} && \psol{$x - \dfrac{x^2}{2}$} && \psol{$(-1, 1]$} \\ \cline{3-3} \cline{5-5} \cline{7-7}
	\\[1.5cm]
	$\arctan(x)$ && \psol{$\ds\sum_{n=0}^\infty (-1)^n \, \dfrac{x^{2n + 1}}{2n + 1}$} && \psol{$x - \dfrac{x^3}{3}$} && \psol{$[-1, 1]$} \\ \cline{3-3} \cline{5-5} \cline{7-7}
	\end{tabular}
	\end{table}



% Question 2
\newpage
\question[16] Compute each of the following: \par\vspace{0.3cm}
	\begin{enumerate}[(a)]
	% (a)
	\item $1 - \dfrac{1}{1!} + \dfrac{1}{2!} - \dfrac{1}{3!} + \dfrac{1}{4!} - \dfrac{1}{5!} + \cdots$ \par\vspace{0.3cm}
		\[
		\begin{gathered}
		\psol{e^x= \sum_{n=0}^\infty \dfrac{x^n}{n!}, \quad x \in \mathbb{R}} \\[0.3cm]
		\psol{e^x= 1 + x + \dfrac{x^2}{2!} + \dfrac{x^3}{3!} + \dfrac{x^4}{4!} + \dfrac{x^5}{5!} + \cdots} \\[0.3cm]
		\psol{e^{-1}= 1 + (-1) + \dfrac{(-1)^2}{2!} + \dfrac{(-1)^3}{3!} + \dfrac{(-1)^4}{4!} + \dfrac{(-1)^5}{5!} + \cdots} \\[0.3cm]
		\psol{e^{-1}= 1 - 1 + \dfrac{1}{2!} - \dfrac{1}{3!} + \dfrac{1}{4!} - \dfrac{1}{5!} + \cdots} \\[0.3cm]
		\psol{\boxed{\dfrac{1}{e}= 1 - 1 + \dfrac{1}{2!} - \dfrac{1}{3!} + \dfrac{1}{4!} - \dfrac{1}{5!} + \cdots}}
		\end{gathered}
		\] \par\vspace{2.25cm}
	
	% (b)
	\item $\ds\sum_{n=1}^\infty \dfrac{(-\pi^2)^n}{(2n)!}$ \par\vspace{0.3cm}
		\[
		\begin{gathered}
		\psol{\cos x= \sum_{n=0}^\infty (-1)^n \, \dfrac{x^{2n}}{(2n)!}, \quad x \in \mathbb{R}} \\
		\psol{\cos x= (-1)^0 \, \dfrac{x^{0}}{0!} + \sum_{n=1}^\infty (-1)^n \, \dfrac{x^{2n}}{(2n)!}} \\
		\psol{\cos x= 1 + \sum_{n=1}^\infty \dfrac{(-1)^n (x^2)^n}{(2n)!}} \\ 
		\psol{\cos x= 1 + \sum_{n=1}^\infty \dfrac{(-x^2)^n}{(2n)!}} \\ 
		\psol{\cos \pi= 1 + \sum_{n=1}^\infty \dfrac{(-\pi^2)^n}{(2n)!}} \\ 
		\psol{-1= 1 + \sum_{n=1}^\infty \dfrac{(-x^2)^n}{(2n)!}} \\
		\psol{\boxed{-2= \sum_{n=1}^\infty \dfrac{(-x^2)^n}{(2n)!}}}
		\end{gathered}
		\]
	
	\newpage
	
	% (c)
	\item $\ds\sum_{n=1}^\infty \dfrac{(-1)^n}{n\, 5^n}$ \par\vspace{0.3cm}
		\[
		\begin{gathered}
		\psol{\ln(1 + x)= \sum_{n=0}^\infty (-1)^{n+1} \, \dfrac{x^n}{n}, \quad x \in (-1, 1]} \\[0.2cm]
		\psol{\ln(1 + x)= (-1) \sum_{n=0}^\infty (-1)^n \, \dfrac{x^n}{n}} \\[0.2cm]
		\psol{-\ln(1 + x)= \sum_{n=0}^\infty (-1)^n \, \dfrac{x^n}{n}} \\[0.2cm]
		\psol{-\ln\left(1 + \tfrac{1}{5} \right)= \sum_{n=0}^\infty (-1)^n \, \dfrac{\left( \frac{1}{5} \right)^n}{n}} \\[0.2cm]
		\psol{-\ln\left(\tfrac{6}{5} \right)= \sum_{n=0}^\infty \dfrac{(-1)^n}{n\, 5^n}} \\[0.2cm]
		\psol{\boxed{\ln\left(\dfrac{5}{6} \right)= \sum_{n=0}^\infty \dfrac{(-1)^n}{n\, 5^n}}}
		\end{gathered}
		\] \par\vspace{0.4cm}
	
	\item $8 - \dfrac{8}{3} + \dfrac{8}{5} - \dfrac{8}{7} + \dfrac{8}{9} - \dfrac{8}{11} + \cdots$ \par\vspace{0.3cm}
		\[
		\begin{gathered}
		\psol{\arctan x= \sum_{n=0}^\infty (-1)^n\, \dfrac{x^{2n+1}}{2n + 1}, \quad x \in [-1, 1]} \\[0.3cm]
		\psol{\arctan x= x - \dfrac{x^3}{3} + \dfrac{x^5}{5} - \dfrac{x^7}{7} + \dfrac{x^9}{9} - \dfrac{x^{11}}{11} + \cdots} \\[0.3cm]
		\psol{8\arctan x= 8x - \dfrac{8x^3}{3} + \dfrac{8x^5}{5} - \dfrac{8x^7}{7} + \dfrac{8x^9}{9} - \dfrac{8x^{11}}{11} + \cdots} \\[0.3cm]
		\psol{8 \arctan(1)= 8(1) - \dfrac{8(1^3)}{3} + \dfrac{8(1^5)}{5} - \dfrac{8(1^7)}{7} + \dfrac{8(1^9)}{9} - \dfrac{8(1^{11})}{11} + \cdots} \\[0.3cm]
		\psol{8 \left( \tfrac{\pi}{4} \right)= 8 - \dfrac{8}{3} + \dfrac{8}{5} - \dfrac{8}{7} + \dfrac{8}{9} - \dfrac{8}{11} + \cdots} \\[0.3cm]
		\psol{\boxed{2\pi = 8 - \dfrac{8}{3} + \dfrac{8}{5} - \dfrac{8}{7} + \dfrac{8}{9} - \dfrac{8}{11} + \cdots}}
		\end{gathered}
		\]
	\end{enumerate}



% Question 3
\newpage
\question[16] Find the center, radius of convergence, and interval of converge for the following power series:
	\[
	\sum_{n=1}^\infty \dfrac{(-1)^n}{n\, 3^n}\; (x + 7)^n
	\] \par\vspace{0.3cm}

\vsol{\itshape Clearly, the center is $x= -5$ --- the $x$-value which `kills' the series. Using the Root Test, we have\dots
	\[
	\lim_{n \to \infty} \left| \dfrac{(-1)^n}{n\, 3^n}\; (x + 7)^n \right|^{1/n}= \lim_{n \to \infty} \left| \dfrac{1}{\sqrt[n]{n} \, 3}\ (x + 7) \right|= \left| \dfrac{1}{1(3)} \, (x + 7) \right|= \left| \dfrac{x + 7}{3} \right|
	\]
Alternatively, using the Root Test, we have\dots
	\[
	\begin{aligned}
	\lim_{n \to \infty} \left| \dfrac{\;\;\dfrac{(-1)^{n+1}}{(n+1)\, 3^{n+1}}\; (x + 7)^{n+1}\;\;}{\dfrac{(-1)^n}{n\, 3^n}\; (x + 7)^n} \right|&= \lim_{n \to \infty} \left| \dfrac{(x + 7)^{n+1}}{(x + 7)^n} \cdot \dfrac{n}{n + 1} \cdot \dfrac{3^n}{3^{n+1}} \right| \\
	&= \lim_{n \to \infty} \left| (x + 7) \cdot \dfrac{n}{n + 1} \cdot \dfrac{1}{3} \right| \\
	&= \left| (x + 7) \cdot 1 \cdot \dfrac{1}{3} \right| \\
	&= \left| \dfrac{x + 7}{3} \right|
	\end{aligned}
	\]
We know that the series converges if this limit is less than one, i.e. $\left| \tfrac{x + 7}{3} \right| < 1$. But then $-1 < \tfrac{x + 7}{3} < 1$, which implies $-3 < x + 7 < 3$. Therefore, the series converges must converge on $-10 < x < -4$. \par\vspace{0.3cm}

We now test the endpoints. If $x= -10$, the series is\dots
	\[
	\sum_{n=1}^\infty \dfrac{(-1)^n}{n\, 3^n}\; (-10 + 7)^n= \sum_{n=1}^\infty \dfrac{(-1)^n}{n\, (-3)^n}\; (-3)^n= \sum_{n=1}^\infty \dfrac{(-1)^n}{n\, 3^n}\; (-1)^n\, 3^n= \sum_{n=1}^\infty \dfrac{1}{n}
	\]
This series diverges by the $p$-test with $p= 1$ (or the fact that the Harmonic series diverge). If $x= -4$, we have\dots
	\[
	\sum_{n=1}^\infty \dfrac{(-1)^n}{n\, 3^n}\; (-4 + 7)^n= \sum_{n=1}^\infty \dfrac{(-1)^n}{n\, 3^n}\; 3^n= \sum_{n=1}^\infty \dfrac{(-1)^n}{n}
	\]
This series is alternating. Observe $\ds\lim_{n \to \infty} \tfrac{1}{n}= 0$ and the sequence $\left\{ \tfrac{1}{n} \right\}$ is decreasing. Therefore, the series converges by the Alternating Series Test. \par\vspace{0.3cm}

So, the interval of convergence is $(-10, -4]$. The radius of convergence is $R= \tfrac{-4 - (-10)}{2}= \tfrac{-4 + 10}{2}= \tfrac{6}{2}= 3$.
	\[
	\boxed{%
	\begin{gathered}
	\text{Center: } x= -7 \\
	\text{Interval of Convergence: } (-10, -4] \\
	\text{Radius of Convergence, } R= 3
	\end{gathered}}
	\]
}



% Question 4
\newpage
\question[16] Use a known Maclaurin series to find a power series representation for each of the expressions below. Be sure to fully simplify your series. \par\vspace{0.2cm}
	\begin{enumerate}[(a)]
	% (a) 
	\item $3x \sin(x^5)$ \par\vspace{0.3cm}
		\[
		\begin{gathered}
		\psol{\sin x= \sum_{n=0}^\infty (-1)^n\, \dfrac{x^{2n+1}}{(2n + 1)!}} \\[0.2cm]
		\psol{\sin(x^5)= \sum_{n=0}^\infty (-1)^n\, \dfrac{(x^5)^{2n+1}}{(2n + 1)!}} \\[0.2cm]
		\psol{\sin(x^5)= \sum_{n=0}^\infty (-1)^n\, \dfrac{x^{10n+5}}{(2n + 1)!}} \\[0.2cm]
		\psol{3x \sin(x^5)= 3x \sum_{n=0}^\infty (-1)^n\, \dfrac{x^{10n+5}}{(2n + 1)!}} \\[0.2cm]
		\psol{\boxed{3x \sin(x^5)= \sum_{n=0}^\infty (-1)^n\, \dfrac{3x^{10n+6}}{(2n + 1)!}}} \\
		\end{gathered}
		\] \par\vspace{1.2cm}
	
	% (b) 
	\item $\dfrac{5}{3 + x}$ \par\vspace{0.2cm}
		\[
		\begin{gathered}
		\psol{\dfrac{5}{3 + x}= \dfrac{5}{3 + x} \cdot \dfrac{1/3}{1/3}} \\[0.2cm]
		\psol{\dfrac{5}{3 + x}= \dfrac{5/3}{1 + \tfrac{x}{3}}} \\[0.2cm]
		\psol{\dfrac{5}{3 + x}= \dfrac{5}{3} \cdot \dfrac{1}{1 - \left( -\tfrac{x}{3} \right)}} \\[0.2cm]
		\psol{\dfrac{5}{3 + x}= \dfrac{5}{3} \sum_{n=0}^\infty \left( -\dfrac{x}{3} \right)^n} \\[0.2cm]
		\psol{\dfrac{5}{3 + x}= \sum_{n=0}^\infty \dfrac{5}{3} \cdot (-1)^n \cdot \dfrac{x^n}{3^n}} \\[0.2cm]
		\psol{\boxed{\dfrac{5}{3 + x}= \sum_{n=0}^\infty \dfrac{5(-1)^n x^n}{3^{n+1}}}}
		\end{gathered}
		\]
	
	\newpage
	
	% (c)
	\item $3x^5 e^{-x/3}$ \par\vspace{0.1cm}
		\[
		\begin{gathered}
		\psol{e^x= \sum_{n=0}^\infty \dfrac{x^n}{n!}} \\[0.2cm]
		\psol{e^{-x/3}= \sum_{n=0}^\infty \dfrac{\left(-\tfrac{x}{3} \right)^n}{n!}} \\
		\psol{e^{-x/3}= \sum_{n=0}^\infty (-1)^n\, \dfrac{x^n}{n! \, 3^n}} \\
		\psol{3x^5 e^{-x/3}= 3x^5 \sum_{n=0}^\infty (-1)^n\, \dfrac{x^n}{n! \, 3^n}} \\
		\psol{3x^5 e^{-x/3}= \sum_{n=0}^\infty (-1)^n\, \dfrac{3x^{n+5}}{n! \, 3^n}} \\
		\psol{\boxed{3x^5 e^{-x/3}= \sum_{n=0}^\infty (-1)^n\, \dfrac{x^{n+5}}{n! \, 3^{n-1}}}} 
		\end{gathered}
		\] \par\vspace{0.3cm}
	
	% (d)
	\item $\ds\int \dfrac{\sin x}{x} \;dx$ \par\vspace{0.1cm}
		\[
		\begin{gathered}
		\psol{\sin x= \sum_{n=0}^\infty (-1)^n\, \dfrac{x^{2n+1}}{(2n + 1)!}} \\
		\psol{\dfrac{\sin x}{x}= \dfrac{1}{x} \,\sum_{n=0}^\infty (-1)^n\, \dfrac{x^{2n+1}}{(2n + 1)!}} \\
		\psol{\dfrac{\sin x}{x}= \sum_{n=0}^\infty (-1)^n\, \dfrac{x^{2n}}{(2n + 1)!}} \\
		\psol{\int \dfrac{\sin x}{x} \;dx= \int \sum_{n=0}^\infty (-1)^n\, \dfrac{x^{2n}}{(2n + 1)!} \;dx} \\
		\psol{\int \dfrac{\sin x}{x} \;dx= \sum_{n=0}^\infty \int (-1)^n\, \dfrac{x^{2n}}{(2n + 1)!} \;dx} \\
		\psol{\int \dfrac{\sin x}{x} \;dx= \sum_{n=0}^\infty \dfrac{(-1)^n}{(2n + 1)!} \int x^{2n} \;dx} \\
		\psol{\int \dfrac{\sin x}{x} \;dx= C + \sum_{n=0}^\infty \dfrac{(-1)^n}{(2n + 1)!} \cdot \dfrac{x^{2n+1}}{2n + 1} \;dx} \\
		\psol{\boxed{\int \dfrac{\sin x}{x} \;dx= C + \sum_{n=0}^\infty (-1)^n\, \dfrac{x^{2n + 1}}{(2n + 1)! \, (2n+1)}}}
		\end{gathered}
		\]
	\end{enumerate}



% Question 5
\newpage
\question[16] Showing all your work, complete the following parts: 
	\begin{enumerate}[(a)]
	% (a)	
	\item Find the first-degree Taylor polynomial, $T_1(x)$, for $\ln x$ centered at $x= 5$. \par\vspace{0.3cm}
		\[
		\begin{aligned}
		&\psol{f(x)= \ln x \bigg|_{x=5}= \ln 5} \\[0.2cm]
		&\psol{f'(x)= \dfrac{1}{x} \bigg|_{x=5}= \dfrac{1}{5}}
		\end{aligned}
		\] \par\vspace{0.2cm}
		\[
		\psol{T_1(x)= \dfrac{f(5)}{0!}\, (x - 5)^0 + \dfrac{f'(5)}{1!}\, (x - 5)^1= \ln 5 + \dfrac{1}{5} \, (x - 5)}
		\]
		\[
		\psol{\boxed{T_1(x)= \ln 4 + \dfrac{1}{4} \, (x - 4)}}
		\] \par\vspace{0.3cm}
	
	% (b)
	\item Use the Taylor Remainder Theorem to find the maximum error in approximating $\ln 6$  by $T_1(6)$. \par\vspace{0.3cm}
	\psol{{\itshape By the Taylor Remainder Theorem, using the fact that $f''(x)= -\frac{1}{x^2}$, \dots}}
		\[
		\psol{R_1(x)= \dfrac{f''(c)}{2!} \,(x - 5)^2= \dfrac{-1/c^2}{2} \,(x - 5)^2= -\dfrac{1}{2c^2} \,(x - 5)^2}
		\]
	\psol{{\itshape where $c$ is some number between the center $4$ and $x$. But then if $x= 6$, $c$ is between $5$ and $6$, and\dots}}
		\[
		\psol{|R_1(5)|= \left| -\dfrac{1}{2c^2} \,(6 - 5)^2 \right|= \dfrac{1}{2c^2} \leq \dfrac{1}{2(5^2)}= \boxed{\dfrac{1}{50} \approx 0.02}}
		\] \par\vspace{0.7cm}
	
	\item Use the Taylor Remainder Theorem to find the maximum error in approximating $\ln x$ with $T_1(x)$ on $[3, 6]$. \par\vspace{0.3cm}
	\psol{{\itshape We know from (b) that $R_1(x)= -\dfrac{1}{2c^2} \,(x - 5)^2$, where $c$ is between the center $4$ and $x$. If $x \in [3, 6]$, then it must be that $c \in [2, 6]$. But then\dots}}
		\[
		\psol{|R_1(x)|= \left| -\dfrac{1}{2c^2} \,(x - 5)^2 \right|= \dfrac{1}{2c^2} \cdot |x - 5|^2 \leq \dfrac{1}{2(3^2)} \cdot |x - 5|^2= \dfrac{1}{2(3^2)} \cdot 2^2= \boxed{\dfrac{2}{9}= 0.222222}}
		\]
	\end{enumerate}



% Question 6
\newpage
\question[16] Find the Taylor $T(x)$ for $\dfrac{1}{(1 - x)^2}$ centered at $x= 0$, and determine its interval of convergence. Show all your work. You may check your using a known Taylor series, but \textit{no credit} will be given for only using the Taylor series of another function, i.e. a ``memorized'' Taylor series. \pspace

\vsol{\itshape We know that the Taylor series for $f(x)$ centered at $x= c$ is\dots
	\[
	T(x)= \sum_{n=0}^\infty \dfrac{f^{(n)}(c)}{n!} \, (x - c)^n
	\]
We only need to find $f^{(n)}(c)$. But\dots
	\[
	\begin{aligned}
	f(x)&= \dfrac{1}{(1 - x)^2} \\[0.2cm]
	f'(x)&= \dfrac{1 \cdot -2}{(1 - x)^3} \cdot -1= \dfrac{1 \cdot 2}{(1 - x)^3} \\[0.2cm]
	f''(x)&= \dfrac{1 \cdot 2 \cdot -3}{(1 - x)^4} \cdot -1= \dfrac{1 \cdot 2 \cdot 3}{(1 - x)^4} \\[0.2cm]
	f'''(x)&= \dfrac{1 \cdot 2 \cdot 3 \cdot -4}{(1 - x)^5} \cdot -1= \dfrac{1 \cdot 2 \cdot 3 \cdot 4}{(1 - x)^5}
	\end{aligned}
	\] \par\vspace{0.2cm}
Therefore, we can see that $f^{(n)}(x)= \dfrac{(n+1)!}{(1 - x)^{n+2}}$, so that $f^{(n)}(0)= \dfrac{(n+1)!}{(1 - 0)^{n+2}}= (n + 1)!$. So, we have\dots
	\[
	T(x)= \sum_{n=0}^\infty \dfrac{f^{(n)}(0)}{n!} \, (x - 0)^n= \sum_{n=0}^\infty \dfrac{(n + 1)!}{n!} \, x^n= \sum_{n=0}^\infty \dfrac{(n + 1)\, n!}{n!} \, x^n= \sum_{n=0}^\infty (n + 1)\, x^n
	\]
Using the Ratio Test, we see\dots
	\[
	\lim_{n \to \infty} \left| \dfrac{\;\;(n + 2) x^{n+1}\;\;}{(n + 1)x^n} \right|= \lim_{n \to \infty} \left| \dfrac{n + 2}{n + 1} \cdot x \right|= |x|
	\]
We know this series converges if $|x| < 1$, i.e. $-1 < x < 1$. If $x= 1$, the series is $\sum_{n=0}^\infty (n + 1)$, which clearly diverges by the Divergence Test. If $x= -1$, the series is $\ds\sum_{n=0}^\infty (n + 1) (-1)^n$, which clearly diverges by the Divergence Test. Therefore, the interval of convergence is $(-1, 1)$, i.e. $-1 < x< 1$. 
	\[
	\boxed{T(x)= \sum_{n=0}^\infty (n + 1)\, x^n, \qquad x \in (-1, 1)}
	\]
}

\end{questions}
\end{document}